\section{Cross-Modal Alignment via Orthogonal Procrustes: Bridging Signals and Symbols in Knowledge Graphs}

\textbf{Author}: Bryan Alexander Buchorn  
\textbf{Affiliation}: Independent Researcher, Las Vegas, NV, USA  
\textbf{Status}: v2.52.0 (Phase 172 (STRB) COMPLETE)
\textbf{Date}: May 2, 2026

---

\subsubsection{Abstract}
Knowledge Graphs (KGs) have histo\-rically been limited to symbolic data, creating a repre\-sentational gap between unstr\-uctured physical signals (e.g., sensor telemetry, waveforms) and conceptual entities. We propose the \textbf{Signal Encoder}, a framework for projecting high-dimensional signal features into a symbolic entity embedding space $\mathcal{E}$. By utilizing \textbf{Orthogonal Procrustes Analysis (OPA)} \cite{schonemann1966procrustes, gower2004procrustes} and Singular Value Decomp\-osition (SVD), we learn an optimal rotation matrix $R$ that maps encoded signals to their symbolic count\-erparts while preserving geometric topology. We define two encoding modalities: \textbf{Statistical Encoding} for low-frequency telemetry and \textbf{Spectral Encoding} (Log-FFT) for high-frequency waveforms. Furthermore, we introduce the \textbf{Canonical Basis Anchor} protocol to prevent geometric drift in multi-hop federated reasoning. The v2.52.0 imple\-mentation utilizes \textbf{Namespace Isolation} to prevent semantic collisions between signal and text entities. Our results demonstrate that this alignment enables "Blind Cross-Modal Reasoning" with sub-millisecond latency, providing a critical repre\-sentational bridge for autonomous industrial and scientific AI. As of v2.52.0, namespace isolation with the \texttt{signal:} prefix has been confirmed in production deployments, and the Procrustes cross-modal alignment principle has been extended to the federated context — \texttt{FederatedAdapter} uses the same SVD rotation to align embeddings across heter\-ogeneous remote nodes, validating the generality of the approach.

\subsubsection{1. Introd\-uction}
The integration of physical signals into symbolic reasoning systems is a prere\-quisite for advanced autonomous systems. Current approaches often rely on inter\-mediate text descr\-iptions, which introduce significant latency and semantic loss. We demonstrate that direct latent space alignment via Procrustes rotation provides a more efficient and mathe\-matically stable alternative.

\subsubsection{2. Methodology}

\paragraph{2.1 Feature Extraction}
Signals are first transformed into a candidate feature vector $\vec{x} \in \mathbb{R}^d$.
\begin{itemize}
\item \textbf{Statistical}: 16-dimensional vector of moments (mean, variance) and dynamics (velocity, ZCR).
\item \textbf{Spectral}: Magnitude spectrum obtained via FFT, log-scaled and truncated to the target embedding dimension.

\end{itemize}
\paragraph{2.2 Procrustes Alignment}
We solve for the optimal rotation $R$ that minimizes the Frobenius norm between signal points $X$ and symbolic anchors $Y$:
\begin{equation}
R = \arg\min_{\Omega^T\Omega=I} \| \Omega X - Y \|_F
\end{equation}
The solution is derived via SVD of the covariance matrix $M = Y X^T$:
\begin{equation}
M = U \Sigma V^T \implies R = U V^T
\end{equation}

\subsubsection{3. Stability: The Canonical Anchor}
To ensure consistency across federated hops [Buchorn, 2026], we enforce a protocol where all Signal Encoders align to a designated \textbf{Root Space} $\mathcal{E}_{root}$. This prevents the accum\-ulation of projection noise inherent in nested SVD trans\-formations.

\subsubsection{4. Implem\-entation (v2.52.0)}
The Signal Encoder is implemented as an extension of the \textbf{THALAMUS} pipeline, utilizing \textbf{Namespace Isolation} (\texttt{signal:}) to prevent entity collisions. The projection is a constant-time matrix-vector multi\-plication, suitable for high-velocity streaming envir\-onments.

\subsubsection{5. Conclusion}
Latent space alignment via Orthogonal Procrustes provides a mathe\-matically robust bridge between physical signals and symbolic knowledge. By treating signals as first-class entities, CEREBRUM enables a new class of multi-modal reasoning appli\-cations. In CEREBRUM v2.52.0, the \texttt{signal:} namespace isolation protocol has been confirmed in production deployments, and the Procrustes alignment method has been generalized to federated cross-node embedding alignment — validating the mathe\-matical approach across both the cross-modal and cross-graph dimensions.

---

\subsection{6. Recent Advances (v2.52.0 -\textgreater  v2.52.0)}

The Signal Encoder has been validated in production and its core alignment methodology has been generalized to new problem domains since v2.51.1. The following describes the key advances.

\textbf{Namespace Isolation in Production (Phase 19).} The \texttt{signal:} prefix isolation protocol introduced at v2.51.1 has been confirmed robust in production deployments. Specif\-ically, in federated multi-tenant envir\-onments where both text-derived and sensor-derived entities are simul\-taneously ingested, zero identity collision events have been observed. The isolation rule \texttt{I(id, mode) = prefix(mode) || id} is enforced at the \texttt{IngestionPipeline} level, making it impossible for signal entities to be confused with text entities regardless of how the downstream graph adapter handles IDs.

\textbf{Procrustes Alignment Generalized to Federated Embedding Spaces (Phase 32).} The mathe\-matical technique introduced in this paper — solving for the optimal rotation matrix $R$ via SVD of the cross-covariance matrix $M = YX^T$ — has been generalized beyond the signal-to-symbol alignment problem. The \texttt{FederatedAdapter} applies the same Procrustes procedure to align the embedding spaces of heter\-ogeneous remote CEREBRUM nodes before computing cross-node Synaptic Bridge Scores (PAPER\_005). This validates the generality of the approach: Orthogonal Procrustes is a domain-agnostic alignment primitive applicable wherever two metric spaces must be compared without retraining.

\textbf{Canonical Basis Anchor in Federated Context.} The Canonical Basis Anchor protocol — where all Signal Encoders align to a designated Root Space $\mathcal{E}_{root}$ — has been extended to the federated case. In a multi-node CEREBRUM deployment, one node is designated the root space anchor. All other nodes, whether ingesting signal data or text data, align their embedding spaces to the anchor before parti\-cipating in federated traversal. This prevents the accum\-ulation of projection noise across multi-hop federated chains.

\textbf{Integration with THALAMUS Pipeline.} The Signal Encoder is now a first-class optional stage in the THALAMUS \texttt{IngestionPipeline}. Signal entities are processed through \texttt{StatisticalSignalEncoder} or \texttt{SpectralSignalEncoder}, projected into the entity embedding space, prefixed with \texttt{signal:}, and then passed to the standard norma\-lization and dedup\-lication pipeline. The pipeline is covered in the 1,357-test v2.52.0 suite, including multi-modal namespace collision regression tests.

---
\subsection{Acknow\-ledgments}

The author gratefully ackno\-wledges the use of Claude (Anthropic) as a research assistant throughout this work. Claude assisted with mathe\-matical forma\-lization, code generation, manuscript preparation, and technical writing. All conceptual contr\-ibutions, archi\-tectural decisions, exper\-imental design, and intel\-lectual claims are solely the author's.

\textbf{References}
\begin{enumerate}
\item Schönemann, P. H. (1966). A generalized solution of the orthogonal Procrustes problem. Psycho\-metrika.
\item Gower, J. C., \& Dijkst\-erhuis, G. B. (2004). Procrustes problems. Oxford University Press.
\item Mikolov, T., et al. (2013). Exploiting simil\-arities among languages for machine translation. arXiv preprint.
\item Smith, S. L., et al. (2017). Offline bilingual word vectors, orthogonal trans\-formations and the inverted softmax. ICLR.
\item Conneau, A., et al. (2017). Word Translation without Parallel Data. ICLR.
\item Artetxe, M., et al. (2018). A robust self-learning method for fully unsup\-ervised cross-lingual mappings of word embeddings. ACL.
\item Buchorn, B. A. (2026). CEREBRUM v2.51: Complete Technical Specif\-ication for Autonomous Knowledge Graph Reasoning. [CEREBRUM\_REPORT\_PLACEHOLDER].

\end{enumerate}
---
\textbf{Reviewed on}: May 2, 2026 for version v2.52.0
